% 【理论扩展】所属：第4章（由 chapters/revised/04-sources.tex \input 引入）
\minitopic{知觉理论的三个问题}知觉哲学至少包含现象问题、形而上学问题与认识问题。现象问题问经验对主体来说如何；形而上学问题问经验与对象、性质和表征的关系；认识问题问经验怎样提供理由或知识。一个理论可在其中一项有优势，却不能自动解决另两项。 例如，幻觉与真知觉主观不可区分支持共同因素现象描述，但不必证明二者形而上学同类；直接实在论说明对象关系，却仍需解释错误经验；教条主义给经验初步确证角色，却可对经验本体保持中立。阅读时应标出作者究竟在哪一层论证。

\minitopic{感觉材料与经验帷幕}感觉材料理论认为，我们直接意识到某种心灵依赖对象，外物知识通过它获得。错觉论证是：水中硬币看起来椭圆；硬币本身仍圆；我们直接意识到的似乎不是硬币实际形状，而是一个椭圆呈现对象。若真知觉与错觉共享直接对象，外物退到经验帷幕后。 反对者说，从“硬币看起来椭圆”不能推出“有一个椭圆对象”。性质看似不同可由观看条件解释。直接实在论允许对象以某种方式呈现，而无需引入内部对象。语法中的“看起来”不自动承诺一个所见的感觉材料。 即使接受内部表征，也不必陷入怀疑。地图能表征城市而不阻止我们借地图知道城市；关键是表征如何可靠连接对象。经验帷幕问题针对的是主体只有表征且无法说明其世界联系的组合，而非任何表征论。

\minitopic{意向论}意向论把知觉经验理解为具有表征内容：经验把世界呈现为某种样子。真知觉与幻觉可具有相同内容，差异在内容是否准确及其成因。它统一解释经验如何支持命题，也容纳不存在对象的幻觉。 困难之一是现象特征能否完全由内容解释。两个经验可能表征同一物体，却因注意、清晰度或感官通道不同而感觉不同。理论可扩充内容、加入表征方式，或承认非意向现象成分。扩充过多则有把所有差异按定义放入内容之嫌。 内容还涉及精细度。经验是否表征“朱红色”，即使主体没有朱红概念？若是，经验内容可非概念；若否，专家训练是否改变看见本身还是只改变判断？这与知觉学习和认知渗透直接相关。

\minitopic{析取论的正面主张}析取论不只是说真知觉和幻觉不同，而主张真知觉的根本性质包含与对象的事实性关系。正常看见苹果时，苹果本身部分构成经验；幻觉只有不可区分的表面。二者不能由一个更基本的共同心理状态穷尽\parencite{martin2004}。 这为知觉知识提供强资源：主体的理由可以是看见事实如此，而非一个同样适合真假世界的中性经验。怀疑者问，主体无法从内部判定自己在哪个析取支，如何合理使用事实性理由？析取论者可说，能拥有理由不要求先证明拥有；正常情形中世界本身可进入理由。 反对意见还包括解释幻觉。若幻觉只由“与真知觉不可区分”消极刻画，理论是否说明其现象？析取论可提供多种幻觉原因，只拒绝这些原因产生与真知觉同类的根本状态。其理论负担是给出统一程度与差异程度的平衡。

\minitopic{认知渗透}若主体的信念、欲望或社会偏见改变经验内容，经验再支持原信念可能形成恶性循环。口渴者把灰色图形看得更像水，警员因刻板印象把无害物误见为武器，都是候选。需区分真正改变经验与只改变注意、记忆或判断。 不是所有自上而下影响都坏。专家知识使医生在影像中看出病灶，语言学习使听者分辨音位。认识评价取决于影响是否对真理敏感、能否被反馈校正。认知渗透不是一个统一击败者，而是一类机制。 若经验被有责偏见塑造，主体可能对后续信念负间接责任；若影响完全自动且不可知，责备减弱，但知识仍可能因不可靠而失败。再次需要区分责任与知识。

\minitopic{知觉学习与专家观看}专家并非只拥有更多背景命题。他们的注意搜索、特征分组和眼动模式都会改变。鸟类学家迅速辨认种类，棋手看见有意义局面，说明能力可通过训练把环境信息组织成可用单位。 专业化也会造成盲点。熟悉典型病征可能忽略罕见异常，期待会使含混影像向既有诊断倾斜。良好训练不仅增加模式库，还培养异常检测、双盲复核和知道何时转交他人。 专家知觉挑战纯粹个人内部模型。教材、标注标准、仪器和同行反馈共同塑造能力。其知识既是个人德性表现，也是社会性认知基础设施的产物。

\minitopic{多感官与行动}知觉通常是多感官的。看到杯子并触摸其边缘，两个来源相互校正；听到声音与看到口型也会整合，有时产生麦格克效应等系统错觉。把视觉孤立为唯一模型，会遗漏跨感官权重。 行动也提供反馈。拿起杯子检验距离，绕到物体背后确认形状，说明知觉不是被动快照。能动理论强调，主体通过掌握感觉变化与行动的关系接触世界。即使不接受完整理论，这种观点也显示主动查验如何提升安全性。 技术界面可能阻断行动反馈。远程摄像头只给固定角度，使用者无法改变视点，证据比现场观察弱。设计可提供缩放、原始信号和多角度，以恢复部分探索能力。

\minitopic{综合案例：医学影像中的偶然病灶}医生在增强图像上看见阴影并诊断肿瘤。实际原始图像没有该阴影，算法因噪声生成伪影；患者却在旁边位置真有一个尚未被图像捕获的肿瘤。因此“患者有肿瘤”为真，医生相信，也有专业经验支持，但真值与所见病灶错位。 这是知觉型Gettier案例。直接因果链不适当，经验内容部分错误，信念在近邻情形不安全。若医生查看原始图、使用另一成像方法并定位真实肿瘤，新的多源证据可产生知识。案例提醒技术增强不是透明窗口。 责任评价取决于界面是否标明生成性处理、行业是否要求原图复核、医生是否受过训练。若厂商隐瞒算法性质，个人无责不等于系统认识可靠；制度应保留变换记录和异常申诉。

\begin{tcolorbox}[breakable,colback=white,colframe=blue!30!black,title={知觉案例分析表},fonttitle=\bfseries]
区分对象、经验内容与最终判断；说明是错觉、幻觉还是误认；检查注意和背景信念影响；记录行动与多感官校正；若有仪器，追踪原始信号和变换；分别评价初步理由、可靠性、安全性与责任。
\end{tcolorbox}
